\subsection{Aufgaben}
\begin{exercise} [Äquivalenzklassen regulärer Hexaeder] \label{ex:äquivalenzklassen_regulärer_hexaeder}
Betrachten Sie einen regulären Hexaeder. Überlegen Sie sich wie viele Äquivalenzklassen der Tetraeder unter der Symmetriegruppe $H$ besitzt.
\end{exercise}

\begin{exercise} [Äquivalenzklassen] \label{ex:äquivalenzklassen}
Überlegen Sie sich ein Beispiel, welches mehrere Äquivalenzklassen besitzt.
\end{exercise}



\begin{exercise}[Symmetrie] \label{ex:notwendigkeitsymmetrie}
Gehen Sie nochmals durch den Beweis von Satz \ref{satz:sum_projection_independent_of_plane}. Können Sie erklären wo die Symmetriegruppeneigenschaft notwendig ist? Können Sie auch ein Beispiel nennen, wo diese Symmetrieeigenschaft nicht vorhanden ist und die Rechnung nicht aufgeht?
\end{exercise}

\begin{exercise} \label{ex:alternativbeweislemma}
Gehen Sie den Alternativbeweis zu Lemma (\ref{lemma:quadriertesummeinvarianzstrecken}) durch. \begin{itemize}
    \item[a)] Erklären Sie wieso aus 
\begin{align}
    M = \lambda I_3 \nonumber
\end{align}
folgt.
\item[b)] Berechnen Sie $\lambda$ für $\Omega=H$.
\end{itemize}
\end{exercise}

\begin{exercise}[Orthogonale Projektion] \label{ex:notwendigkeitorthogonaleprojektionen}
Gehen Sie nochmals durch den Beweis von Satz \ref{satz:sum_projection_independent_of_plane}. Erkläre in eigenen Worten, wo die orthogonal Projektionen verwendet werden im Beweis.
\end{exercise}
\begin{exercise}[Umkehrung] \label{ex:umkehrung_satz}
Gilt die Umkehrung von Satz \ref{satz:sum_projection_independent_of_plane}, d.h. falls Gleichung (\ref{eq:sum_projection_independent_of_plane}) gilt, folgt, dass $P \subset \mathbb{R}^3$ ein Polyeder mit Symmetriegruppe $T,H$ oder $I$ ist und wir eine orthogonale Abbildung haben.
\end{exercise}


\begin{exercise}  \label{ex:summe_einheitswurzel}
Zeige, dass für  $n \geq 3$ 

    \begin{align} \label{eq:summe_einheitswurzel}
        \sum_{k=0}^{n-1}e^{i4k\pi/n} = 0  
    \end{align}
gilt.
\end{exercise}


\begin{exercise} \label{ex:summ_einheitswurzel_smaller_than_3}
    Zeige, dass Gleichung (\ref{eq:summe_einheitswurzel}) für $n=1$ und $n=2$ nicht gilt.
\end{exercise}

\begin{exercise} \label{ex:PseudoRhombenkuboktaeder_Symmetriegruppe}
    Zeige, dass das Pseudo-Rhombenkuboktaeder weder Symmetriegruppe H noch I hat.
\end{exercise}